\section{Discussion and Limitations}
\label{sec:discussion}

\subsection{Computational cost}
The full HAVI-Methyl architecture is expected to be more expensive per training step than a per-sample HMM because it runs a fragment-bag encoder and flow posterior. The compensating advantage is amortization: a new sample can be processed by forward passes rather than a full per-sample Baum-Welch fit. Exact ratios, wall-clock times, and accelerator throughput remain implementation-dependent planning estimates until the full architecture is benchmarked. The wafer-scale recipe of Section~\ref{sec:benchmark} should therefore be treated as an implementation plan, not as a completed performance claim.

\subsection{Failure modes}
Several regimes warrant caution. Deeply hypo-methylated partially methylated domains can saturate fragmentomic signals and should be flagged from public PMD tracks before downstream clock or tissue-of-origin analysis. Centromeric and pericentromeric repeats suffer from collapsed mappability, detectable from per-locus mappability scores. Short reads below the chosen QC threshold carry degraded motif information. Low-coverage tail loci with $n^{\mathrm{frag}}<2$ should be reported as posterior-prior fallbacks rather than informative posteriors. The implementation should emit diagnostics for these cases: PMD overlap, mappability flags, read-length distributions, local coverage, posterior effective sample size, and stratum-specific calibration.

\subsection{Extensions}
Several extensions are natural follow-ups. Nanopore direct-methylation likelihoods can replace the Beta-Binomial pseudo-likelihood with the Bernoulli direct-call form. Joint inference with ichorCNA \citep{adalsteinsson2017ichor} could provide CNV-corrected methylation values for cancer cfDNA. A methylation foundation model such as CpGPT or MethylGPT can be attached downstream as an auxiliary consumer of the HAVI-Methyl posterior $\beta$ output, rather than as a direct WGS-fragment baseline.

\subsection{Clinical translation}
Clinical translation of an ML-based methylation predictor requires prospective validation, subgroup analysis, calibration monitoring, and regulatory review. The conformal wrapper supplies a distribution-free marginal coverage guarantee under exchangeability, but deployment still requires empirical auditing across age, sex, ancestry, cohort, batch, and disease strata. The mQTL anchor loss is only as representative as the genotype panels and summary statistics supporting it, so ancestry-specific or multi-ancestry anchor evaluations should be planned. Regulatory details should be verified against current guidance before final claims are made.

\subsection{Current-development limitation}
The manuscript currently contains a full proposed method, proof sketches and conditional theory, completed simplified synthetic experiments, and a planned real-data benchmark. Full architecture implementation, real-data training, conformal calibration evaluation, tissue-of-origin head validation, and all publication-ready figures/tables remain future work for the next iteration.

% ============================================================
